%% ============================================================
%%  DOCUMENT DE CADRAGE LOGIQUE — AGENT HYPOTHÈSES
%%  Version PFE synthétique et professionnelle
%%  3 Chapitres : Introduction & Contexte | Notions Fondamentales | Veille Technologique & Conclusion
%% ============================================================
\documentclass[12pt, a4paper]{report}

%% ---------- Encodage & Langue ----------
\usepackage[utf8]{inputenc}
\usepackage[T1]{fontenc}
\usepackage[french]{babel}

%% ---------- Mise en page ----------
\usepackage[top=2.5cm, bottom=2.5cm, left=2.8cm, right=2.8cm]{geometry}
\usepackage{parskip}
\setlength{\parindent}{0pt}
\setlength{\parskip}{6pt}

%% ---------- Couleurs ----------
\usepackage[dvipsnames,svgnames,table]{xcolor}
\definecolor{navyblue}{HTML}{1A3A5C}
\definecolor{midblue}{HTML}{2E6DA4}
\definecolor{lightblue}{HTML}{D6E8F7}
\definecolor{paleblue}{HTML}{EBF4FB}
\definecolor{darkgreen}{HTML}{1B5E3B}
\definecolor{midgreen}{HTML}{2E7D52}
\definecolor{lightgreen}{HTML}{D4EDDA}
\definecolor{codebg}{HTML}{F3F6FA}
\definecolor{stepbg}{HTML}{EBF4FB}
\definecolor{rulegray}{HTML}{B0BEC5}

%% ---------- Titres personnalisés ----------
\usepackage{titlesec}
\titleformat{\chapter}[display]
  {\normalfont\huge\bfseries\color{navyblue}}
  {\color{midblue}\chaptertitlename\ \thechapter}
  {10pt}
  {\color{navyblue}}
  [\vspace{4pt}\textcolor{midblue}{\hrule height 1.5pt}]

\titleformat{\section}
  {\normalfont\Large\bfseries\color{midblue}}
  {\thesection}{0.8em}{}
  [\vspace{2pt}\textcolor{lightblue}{\hrule height 0.8pt}]

\titleformat{\subsection}
  {\normalfont\large\bfseries\color{darkgreen}}
  {\thesubsection}{0.7em}{}

\titleformat{\subsubsection}
  {\normalfont\normalsize\bfseries\color{midgreen}}
  {\thesubsubsection}{0.6em}{}

%% ---------- En-têtes / Pieds de page ----------
\usepackage{fancyhdr}
\pagestyle{fancy}
\fancyhf{}
\fancyhead[L]{\small\textcolor{midblue}{Agent Hypothèses — Document de Cadrage}}
\fancyhead[R]{\small\textcolor{midblue}{\leftmark}}
\fancyfoot[C]{\textcolor{navyblue}{\thepage}}
\renewcommand{\headrulewidth}{0.4pt}
\renewcommand{\headrule}{\hbox to\headwidth{\color{midblue}\leaders\hrule height \headrulewidth\hfill}}

%% ---------- Tableaux ----------
\usepackage{booktabs}
\usepackage{tabularx}
\usepackage{longtable}
\usepackage{array}
\usepackage{multirow}
\usepackage{colortbl}

\newcolumntype{L}[1]{>{\raggedright\arraybackslash}p{#1}}
\newcolumntype{C}[1]{>{\centering\arraybackslash}p{#1}}

%% ---------- Boîtes colorées ----------
\usepackage{tcolorbox}
\tcbuselibrary{skins, breakable, listingsutf8}

%% Boîte formule / règle
\newtcolorbox{formulebox}[1][]{
  colback=paleblue, colframe=midblue, 
  fonttitle=\bfseries\color{navyblue},
  boxrule=0.8pt, arc=4pt,
  left=8pt, right=8pt, top=6pt, bottom=6pt,
  #1
}

%% Boîte étape (fichier expliqué)
\newtcolorbox{stepbox}[2][]{
  colback=stepbg, colframe=navyblue,
  fonttitle=\bfseries\color{white},
  coltitle=white, attach boxed title to top left={yshift=-2mm, xshift=6mm},
  boxed title style={colback=navyblue, colframe=navyblue, arc=3pt},
  title={#2},
  boxrule=0.6pt, arc=4pt,
  left=8pt, right=8pt, top=14pt, bottom=6pt,
  breakable,
  #1
}

%% Boîte résultat / note verte
\newtcolorbox{resultbox}[1][]{
  colback=lightgreen, colframe=darkgreen,
  boxrule=0.6pt, arc=4pt,
  left=8pt, right=8pt, top=6pt, bottom=6pt,
  #1
}

%% Boîte code
\newtcolorbox{codebox}[1][]{
  colback=codebg, colframe=rulegray,
  boxrule=0.5pt, arc=3pt,
  left=6pt, right=6pt, top=4pt, bottom=4pt,
  fontupper=\small\ttfamily,
  #1
}

%% ---------- Listes ----------
\usepackage{enumitem}
\setlist[itemize,1]{label=\textcolor{midblue}{\textbullet}, leftmargin=1.5em, itemsep=2pt}
\setlist[itemize,2]{label=\textcolor{darkgreen}{$\circ$}, leftmargin=1.5em, itemsep=1pt}
\setlist[enumerate,1]{label=\textcolor{navyblue}{\arabic*.}, leftmargin=1.5em, itemsep=3pt}

%% ---------- Liens ----------
\usepackage[hidelinks]{hyperref}

%% ---------- Divers ----------
\usepackage{graphicx}
\usepackage{amsmath}
\usepackage{amssymb}
\usepackage{setspace}
\usepackage{float}

%% ---------- Table des matières ----------
\usepackage{tocloft}
\renewcommand{\cfttoctitlefont}{\Large\bfseries\color{navyblue}}
\renewcommand{\cftchapfont}{\bfseries\color{navyblue}}
\renewcommand{\cftsecfont}{\color{midblue}}
\renewcommand{\cftsubsecfont}{\color{darkgreen}}
\setlength{\cftbeforechapskip}{6pt}

%% ---------- Arbre de fichiers (verbatim coloré) ----------
\usepackage{fancyvrb}

%% ============================================================
\begin{document}
\onehalfspacing

%% ============================================================
%%  PAGE DE GARDE
%% ============================================================
\begin{titlepage}
\pagecolor{navyblue}
\color{white}
\begin{center}

\vspace*{1.5cm}

{\small\bfseries\color{lightblue} UNIVERSITÉ ABDELMALEK ESSAÂDI — FSTT}\\[4pt]
{\small\color{lightblue} Licence Data Analytics}\\[6pt]
{\small\color{lightblue} Projet de Fin d'Études}

\vspace{1.2cm}
\textcolor{midblue}{\hrule height 1pt}
\vspace{0.3cm}
\textcolor{rulegray}{\hrule height 0.5pt}
\vspace{1.2cm}

{\fontsize{13}{18}\selectfont\bfseries\color{lightblue}
DOCUMENT DE CADRAGE LOGIQUE}\\[10pt]

{\fontsize{26}{34}\selectfont\bfseries\color{white}
Agent Hypothèses}\\[14pt]

{\fontsize{14}{20}\selectfont\color{lightblue}
Module Pédagogique Intelligent}\\[6pt]

{\fontsize{11}{16}\selectfont\color{rulegray}
Plateforme IA Multi-Agents pour l'Accompagnement à la Création d'Entreprise}

\vspace{1.2cm}
\textcolor{rulegray}{\hrule height 0.5pt}
\vspace{0.3cm}
\textcolor{midblue}{\hrule height 1pt}
\vspace{1.8cm}

%% Trois blocs thématiques
\begin{minipage}{0.3\textwidth}
\centering
\textcolor{midblue}{\large\bfseries Chapitre 1}\\[4pt]
\textcolor{lightblue}{\small Introduction \&\\Contexte Stratégique}
\end{minipage}
\hfill
\begin{minipage}{0.3\textwidth}
\centering
\textcolor{midblue}{\large\bfseries Chapitre 2}\\[4pt]
\textcolor{lightblue}{\small Notions Fondamentales\\\& Logique Métier}
\end{minipage}
\hfill
\begin{minipage}{0.3\textwidth}
\centering
\textcolor{midblue}{\large\bfseries Chapitre 3}\\[4pt]
\textcolor{lightblue}{\small Veille Technologique\\\& Conclusion}
\end{minipage}

\vfill

{\small\color{rulegray} Année Universitaire 2025--2026}

\end{center}
\end{titlepage}

\nopagecolor
\color{black}

%% ============================================================
%%  TABLE DES MATIÈRES
%% ============================================================
\tableofcontents
\newpage

%% ============================================================
%%  CHAPITRE 1 — INTRODUCTION & CONTEXTE STRATÉGIQUE
%% ============================================================
\chapter{Introduction \& Contexte Stratégique}

\section{Présentation générale}

L'Agent Hypothèses constitue le cœur intelligent de la plateforme IA multi-agents d'accompagnement à la création d'entreprise au Maroc. Développé dans le cadre d'un Projet de Fin d'Études en Licence Data Analytics, ce module orchestre un dialogue pédagogique structuré pour collecter, valider et transformer les hypothèses financières des entrepreneurs en projections prévisionnelles robustes.

La plateforme repose sur une architecture multi-agents orchestrée par \textbf{n8n}, où chaque agent couvre une phase spécifique du parcours entrepreneurial : idéation, analyse de marché, collecte d'hypothèses, génération financière, conformité juridique et production du Business Plan final. L'Agent Hypothèses représente le pont stratégique entre le dialogue avec l'entrepreneur et la génération des tableaux financiers prévisionnels.

\section{Problématique}

Les entrepreneurs marocains font face à plusieurs défis majeurs dans l'élaboration de leurs prévisions financières :

\begin{itemize}
  \item \textbf{Difficulté à produire des hypothèses réalistes} : manque d'expérience financière et de repères sectoriels
  \item \textbf{Absence d'accompagnement pédagogique} : outils complexes sans guidance adaptée
  \item \textbf{Risque d'erreurs dans les prévisions} : conséquences graves sur la viabilité du projet
  \item \textbf{Besoin d'automatisation intelligente} : nécessité de combiner expertise humaine et efficacité technologique
\end{itemize}

\section{Objectifs}

L'Agent Hypothèses vise à répondre à ces défis à travers plusieurs objectifs précis :

\begin{itemize}
  \item \textbf{Collecte structurée} : recueillir les 22 hypothèses financières essentielles selon la méthode Chauvin
  \item \textbf{Validation métier} : appliquer les règles financières pour garantir la cohérence des prévisions
  \item \textbf{Génération automatique} : produire 3 scénarios financiers (pessimiste, réaliste, optimiste)
  \item \textbf{Production structurée} : générer un JSON normalisé pour l'Agent Finance
  \item \textbf{Intégration fluide} : s'intégrer parfaitement dans l'architecture multi-agents globale
\end{itemize}

\section{Architecture globale du workflow}

\begin{formulebox}[title={\bfseries\color{navyblue} Workflow de la plateforme}]
\texttt{Angular $\rightarrow$ FastAPI $\rightarrow$ SLM Phi-3 Mini (classification intention)}\\
\texttt{$\rightarrow$ Agent Contrôleur n8n $\rightarrow$ POST /hypotheses/start}\\
\texttt{$\rightarrow$ Dialogue 22 questions (RAG Chauvin + Groq)}\\
\texttt{$\rightarrow$ Validation Chauvin $\rightarrow$ JSON HypothesesOutput}\\
\texttt{$\rightarrow$ Agent Finance (POST /plan) $\rightarrow$ Agent Business Plan $\rightarrow$ PDF final}
\end{formulebox}

\section{Arborescence du projet}

\begin{codebox}
agent\_hypotheses/\\
\quad{}+-- .env\\
\quad{}+-- main.py\\
\quad{}+-- requirements.txt\\
\quad{}+-- agent/\\
\quad{}|\quad{}+-- \_\_init\_\_.py\\
\quad{}|\quad{}+-- dialogue.py\\
\quad{}|\quad{}+-- groq\_client.py\\
\quad{}|\quad{}+-- validation.py\\
\quad{}|\quad{}+-- questions.py\\
\quad{}|\quad{}+-- scenarios.py\\
\quad{}|\quad{}\textbackslash{}-- output\_builder.py\\
\quad{}+-- rag/\\
\quad{}|\quad{}+-- \_\_init\_\_.py\\
\quad{}|\quad{}+-- vectorstore.py\\
\quad{}|\quad{}+-- embedder.py\\
\quad{}|\quad{}\textbackslash{}-- retriever.py\\
\quad{}+-- schemas/\\
\quad{}|\quad{}\textbackslash{}-- output\_schema.py\\
\quad{}+-- data/\\
\quad{}|\quad{}\textbackslash{}-- scenarios\_reference.json\\
\quad{}\textbackslash{}-- scripts/\\
\quad{}\quad{}\quad{}+-- test\_dialogue.py\\
\quad{}\quad{}\quad{}+-- test\_validation.py\\
\quad{}\quad{}\quad{}+-- test\_json\_final.py\\
\quad{}\quad{}\quad{}+-- test\_scenarios\_rag.py\\
\quad{}\quad{}\quad{}+-- build\_scenarios\_rag.py
\end{codebox}

\section{Explication des fichiers}

\begin{itemize}
  \item \textbf{main.py} : API FastAPI principale, point d'entrée HTTP avec 5 endpoints
  \item \textbf{dialogue.py} : gestion du dialogue conversationnel et orchestration des 22 questions
  \item \textbf{groq\_client.py} : connexion au LLM Llama-3.3-70b via API Groq
  \item \textbf{validation.py} : validation métier des hypothèses selon les règles Chauvin
  \item \textbf{questions.py} : dictionnaire des 22 questions structurées en 4 blocs
  \item \textbf{scenarios.py} : calcul des 3 projections financières (pessimiste, réaliste, optimiste)
  \item \textbf{output\_builder.py} : construction du JSON final structuré pour l'Agent Finance
  \item \textbf{vectorstore.py} : gestion de la base ChromaDB avec 280 documents indexés
  \item \textbf{embedder.py} : modèle d'embeddings all-MiniLM-L6-v2 (384 dimensions)
  \item \textbf{retriever.py} : recherche sémantique et récupération de contexte pertinent
  \item \textbf{output\_schema.py} : schémas Pydantic définissant la structure JSON exacte
\end{itemize}

\section{Étapes de réalisation}

\begin{stepbox}{Étape 1 — Configuration de l'environnement}
Création de l'environnement virtuel Python 3.10, installation des dépendances FastAPI, Pydantic, Groq, ChromaDB, sentence-transformers. Configuration des variables sensibles dans .env (API keys, chemins bases).
\end{stepbox}

\begin{stepbox}{Étape 2 — Schémas Pydantic}
Définition du contrat JSON formel avec les classes Bloc1Ventes, Bloc2Achats, Bloc3Charges, Bloc4Encaissements, ResultatValidation, Scenario, TroisScenarios et HypothesesOutput. Garantit la cohérence et le typage des données échangées.
\end{stepbox}

\begin{stepbox}{Étape 3 — Pipeline RAG}
Indexation de 280 documents (270 chunks du livre Chauvin + 10 scénarios marocains) dans ChromaDB. Configuration du modèle d'embeddings all-MiniLM-L6-v2 pour la recherche sémantique par similarité cosinus.
\end{stepbox}

\begin{stepbox}{Étape 4 — Validation métier}
Implémentation des 3 règles de Chauvin : seuil de rentabilité, symétrie financement, cohérence prix/fabrication. Définition des bornes min/max pour chaque hypothèse numérique avec messages d'alerte pédagogiques.
\end{stepbox}

\begin{stepbox}{Étape 5 — Scénarios financiers}
Calcul automatique des 3 projections en appliquant les coefficients 0.7/1.0/1.3. Formule de croissance composée pour le CA mensuel, détermination du mois de rentabilité jusqu'à 24 mois.
\end{stepbox}

\begin{stepbox}{Étape 6 — Dialogue intelligent}
Orchestration des 22 questions avec détection d'intention (réponse, question, hors-sujet, vide). Intégration RAG pour enrichir chaque échange, gestion des tentatives maximales (3 par question).
\end{stepbox}

\begin{stepbox}{Étape 7 — API FastAPI}
Exposition des endpoints /health, /hypotheses/start, /hypotheses/validate, /hypotheses/questions et /scenarios/\{id\}. Initialisation unique des clients Groq, ChromaDB et embedder pour optimiser les performances.
\end{stepbox}

\begin{stepbox}{Étape 8 — Tests et validation}
Scripts de test unitaires pour validation, dialogue, JSON final et recherche RAG. Tests de non-régression et validation des métriques de performance (temps de réponse Groq : 1070ms moyen).
\end{stepbox}

\vspace{12pt}
\begin{resultbox}
\textbf{\color{darkgreen} Architecture en 4 couches :}\\
\textbf{API} (\texttt{main.py}) : point d'entrée HTTP, orchestre les appels aux modules.\\
\textbf{Agent} (\texttt{agent/}) : dialogue, validation métier et construction du modèle de sortie.\\
\textbf{RAG} (\texttt{rag/}) : recherche de contexte dans base vectorielle pour enrichir les prompts.\\
\textbf{Schémas} (\texttt{schemas/}) : description formelle du format JSON attendu.
\end{resultbox}

%% ============================================================
%%  CHAPITRE 2 — NOTIONS FONDAMENTALES & LOGIQUE MÉTIER
%% ============================================================
\chapter{Notions Fondamentales \& Logique Métier}

\section{LLM — Large Language Model}

Le LLM \textbf{Llama-3.3-70b-versatile}, accessible via l'API Groq, constitue le cerveau conversationnel de l'agent. Il remplit cinq fonctions essentielles : générer les questions pédagogiques en français naturel, expliquer chaque hypothèse à l'entrepreneur, valider les réponses contextuelles, répondre aux questions posées pendant le dialogue, et détecter l'intention derrière chaque réponse.

\section{RAG — Retrieval-Augmented Generation}

Le pipeline RAG combine recherche sémantique dans ChromaDB et génération de texte par le LLM. Trois étapes : conversion du contexte en vecteur via all-MiniLM-L6-v2, recherche par similarité cosinus parmi 280 documents, injection des passages pertinents dans le prompt Groq pour générer des réponses ancrées dans les sources Chauvin et scénarios marocains.

\section{Embeddings et Similarité Cosinus}

Le modèle \textbf{all-MiniLM-L6-v2} transforme chaque texte en vecteur de 384 dimensions. La similarité cosinus mesure l'angle entre vecteurs — plus l'angle est petit, plus les textes sont sémantiquement similaires. Permet de retrouver des contenus pertinents même avec des mots différents.

\section{ChromaDB — Base Vectorielle}

ChromaDB stocke les embeddings et permet leur recherche par similarité sémantique. Contient 280 documents : 270 chunks du livre Chauvin et 10 scénarios marocains de référence. Diffère de PostgreSQL par sa capacité à rechercher par pertinence sémantique plutôt que par correspondance exacte.

\section{Les 4 Blocs d'Hypothèses}

\begin{center}
\begin{tabularx}{0.95\textwidth}{
  >{\centering\arraybackslash\bfseries\color{navyblue}}C{1.2cm}
  >{\centering\arraybackslash}C{2.2cm}
  X
  >{\centering\arraybackslash}C{2cm}
}
\toprule
\rowcolor{navyblue}\textcolor{white}{\textbf{Bloc}} & \textcolor{white}{\textbf{Question Chauvin}} & \textcolor{white}{\textbf{Contenu}} & \textcolor{white}{\textbf{Hypothèses}} \\
\midrule
\rowcolor{paleblue} 1 & Quoi ? & Ventes, clients, prix, croissance, fidélisation & H1 à H7 \\
2 & Comment ? & Type activité, fabrication, infrastructure, fournisseurs & H8 à H12 \\
\rowcolor{paleblue} 3 & Combien ? & Charges fixes, loyer, salaires, marketing, investissements & H13 à H21 \\
4 & Quand ? & Nature clients, délais encaissement & H22 \\
\bottomrule
\end{tabularx}
\end{center}

\section{Règles de Validation de Chauvin}

\subsection{Règle 1 — Seuil de Rentabilité}

Le seuil de rentabilité est le niveau de ventes minimum pour que l'entreprise couvre exactement ses charges fixes sans réaliser de bénéfice ni de perte.

\begin{formulebox}[title={\bfseries\color{navyblue} Formule du seuil de rentabilité}]
\[
\text{Charges fixes} = H_{13} + H_{14} + H_{15} + H_{16} + H_{17} + \frac{H_{18}}{12}
\]
\[
\text{Marge unitaire} = H_2 \,(\text{prix de vente}) - H_9 \,(\text{coût de fabrication})
\]
\[
\text{Seuil clients minimum} = \frac{\text{Charges fixes}}{\text{Marge unitaire}}
\]
\textbf{Si clients prévus $<$ seuil} $\Rightarrow$ \textcolor{darkgreen}{\textbf{AVERTISSEMENT pédagogique}}
\end{formulebox}

\subsection{Règle 2 — Symétrie du Financement}

Selon Chauvin, si les besoins de financement dépassent le double des apports personnels de l'entrepreneur, la structure financière est déséquilibrée.

\begin{formulebox}[title={\bfseries\color{navyblue} Formule de la symétrie}]
\[
\text{Apports propres} = H_{19} - H_{21}
\]
\textbf{Si } $H_{19} > 2 \times \text{apports propres}$ $\Rightarrow$ \textcolor{darkgreen}{\textbf{AVERTISSEMENT Chauvin}}
\end{formulebox}

\subsection{Règle 3 — Cohérence Prix / Coût de Fabrication}

Si le coût de fabrication unitaire représente plus de 80~\% du prix de vente, la marge brute est insuffisante pour couvrir les charges fixes. L'agent alerte l'entrepreneur et suggère d'augmenter le prix ou de réduire les coûts.

\section{Les 3 Scénarios Financiers}

Conformément à la méthode Chauvin, l'agent génère automatiquement trois versions du prévisionnel financier en appliquant des coefficients multiplicateurs aux hypothèses de base.

\begin{center}
\begin{tabularx}{0.95\textwidth}{
  >{\bfseries\color{navyblue}}L{2.5cm}
  >{\centering\arraybackslash}C{2cm}
  X
  L{4cm}
}
\toprule
\rowcolor{navyblue}\textcolor{white}{\textbf{Scénario}} & \textcolor{white}{\textbf{Coefficient}} & \textcolor{white}{\textbf{Logique}} & \textcolor{white}{\textbf{Usage}} \\
\midrule
\rowcolor{paleblue} Pessimiste & $\times\, 0{,}7$ & 30~\% moins de clients & Vérifier la survie \\
Réaliste & $\times\, 1{,}0$ & Hypothèses telles que saisies & Prévisions de référence \\
\rowcolor{paleblue} Optimiste & $\times\, 1{,}3$ & 30~\% plus de clients & Capacité à croître \\
\bottomrule
\end{tabularx}
\end{center}

Pour chaque scénario, l'agent calcule : le CA mensuel sur 12~mois (croissance composée), le CA annuel total, les charges annuelles, la marge brute, le résultat net et le mois de rentabilité.

\section{Guardrails — Gestion des Intentions}

Les guardrails sont des mécanismes de protection qui contrôlent le flux du dialogue et empêchent l'agent de se bloquer face à des réponses inattendues. L'agent détecte 4 types d'intention :

\begin{center}
\begin{tabularx}{0.95\textwidth}{
  >{\bfseries\color{navyblue}}L{2.5cm}
  X
  L{4.5cm}
}
\toprule
\rowcolor{navyblue}\textcolor{white}{\textbf{Intention}} & \textcolor{white}{\textbf{Description}} & \textcolor{white}{\textbf{Action de l'agent}} \\
\midrule
\rowcolor{paleblue} Réponse & L'entrepreneur répond directement & Valider et enregistrer \\
Question & L'entrepreneur pose une question & Répondre via RAG puis reposer la question \\
\rowcolor{paleblue} Hors sujet & Réponse sans rapport & Signaler et reposer (max 3 tentatives) \\
Vide & Entrée vide ou espaces & Rappeler poliment et reposer \\
\bottomrule
\end{tabularx}
\end{center}

\section{Délégation vers les Autres Agents}

Quand l'entrepreneur pose une question hors du périmètre du livre Chauvin, l'agent détecte la catégorie et indique quel agent spécialisé devrait répondre.

\begin{center}
\begin{tabularx}{0.95\textwidth}{
  >{\ttfamily\small\color{navyblue}}L{2.2cm}
  X
  >{\bfseries}L{3.2cm}
}
\toprule
\rowcolor{navyblue}\textcolor{white}{\textbf{Catégorie}} & \textcolor{white}{\textbf{Exemple}} & \textcolor{white}{\textbf{Agent délégué}} \\
\midrule
\rowcolor{paleblue} chauvin & C'est quoi le seuil de rentabilité~? & Agent Hypothèses \\
general & C'est quoi une SARL~? & Agent Général \\
\rowcolor{paleblue} marche & Mon secteur est-il saturé~? & Agent Marché \\
finance & Comment calculer mon BFR~? & Agent Finance \\
\rowcolor{paleblue} hors\_sujet & Quel temps fait-il à Tanger~? & Refus poli \\
\bottomrule
\end{tabularx}
\end{center}

\section{Résultats obtenus — Validation du système}

\subsection{Métriques de performance}

\begin{center}
\begin{tabularx}{0.95\textwidth}{X >{\centering\arraybackslash}C{3cm} >{\centering\arraybackslash}C{3cm}}
\toprule
\rowcolor{navyblue}\textcolor{white}{\textbf{Indicateur}} & \textcolor{white}{\textbf{Valeur obtenue}} & \textcolor{white}{\textbf{Statut}} \\
\midrule
\rowcolor{paleblue} Connexion Groq & 1~070~ms & \checkmark{} $<$ 2~s \\
Documents RAG & 280 indexés & \checkmark{} \\
\rowcolor{paleblue} Questions couvertes & 22 / 22 & \checkmark{} \\
Endpoints API & 5 / 5 & \checkmark{} \\
\rowcolor{paleblue} Scénarios générés & 3 automatiques & \checkmark{} \\
JSON final & 12 champs validés & \checkmark{} \\
\bottomrule
\end{tabularx}
\end{center}

\subsection{Exemple de résultat — Scénario épicerie Tanger}

Hypothèses saisies : H2 = 150~MAD (prix), H4 = 30~clients, H9 = 80~MAD (coût fabrication), H13 = 3~000~MAD (loyer), H14 = 6~000~MAD (salaires).

\begin{center}
\begin{tabularx}{0.95\textwidth}{X >{\centering\arraybackslash}C{2.8cm} >{\centering\arraybackslash}C{2.8cm} >{\centering\arraybackslash}C{2.8cm}}
\toprule
\rowcolor{navyblue}\textcolor{white}{\textbf{Indicateur}} & \textcolor{white}{\textbf{Pessimiste}} & \textcolor{white}{\textbf{Réaliste}} & \textcolor{white}{\textbf{Optimiste}} \\
\midrule
\rowcolor{paleblue} Clients mois~1 & 21 & 30 & 39 \\
CA mois~1 & 3~150~MAD & 4~500~MAD & 5~850~MAD \\
\rowcolor{paleblue} CA annuel & 67~360~MAD & 96~229~MAD & 125~098~MAD \\
Charges annuelles & \multicolumn{3}{c}{143~900~MAD} \\
\rowcolor{paleblue} Résultat net an~1 & $-$112~465~MAD & $-$98~993~MAD & $-$85~521~MAD \\
Mois rentabilité & Mois~22 & Mois~19 & Mois~16 \\
\bottomrule
\end{tabularx}
\end{center}

\begin{resultbox}
\textbf{\color{darkgreen} Message pédagogique généré~:} Votre projet nécessite 19~mois avant rentabilité dans le scénario réaliste. Assurez-vous d'avoir une trésorerie suffisante pour couvrir les 19~premiers mois d'activité.
\end{resultbox}

%% ============================================================
%%  CHAPITRE 3 — VEILLE TECHNOLOGIQUE & CONCLUSION
%% ============================================================
\chapter{Veille Technologique \& Conclusion}

\section{Stack Technique}

\begin{center}
\begin{tabularx}{0.95\textwidth}{
  >{\bfseries\color{navyblue}}L{2.5cm}
  >{\centering\arraybackslash}C{2cm}
  X
  >{\centering\arraybackslash}C{2.5cm}
}
\toprule
\rowcolor{navyblue}\textcolor{white}{\textbf{Composant}} & \textcolor{white}{\textbf{Technologie}} & \textcolor{white}{\textbf{Justification}} & \textcolor{white}{\textbf{Alternative}} \\
\midrule
\rowcolor{paleblue} API REST & FastAPI & Performance, typage automatique & Flask, Django \\
LLM & Llama-3.3-70b & Qualité français, vitesse via Groq & GPT-4, Claude \\
\rowcolor{paleblue} Base vectorielle & ChromaDB & Open-source, Python natif & Pinecone, Weaviate \\
Embeddings & all-MiniLM-L6-v2 & Léger (80Mo), performant & text-embedding-ada-002 \\
\rowcolor{paleblue} Validation & Pydantic & Contrat JSON strict & Marshmallow, Cerberus \\
Orchestration & n8n & Visuel, multi-agents & Airflow, Prefect \\
\bottomrule
\end{tabularx}
\end{center}

\section{Comparaison d'outils}

\subsection{LLM : Groq vs Alternatives}

\begin{itemize}
  \item \textbf{Groq + Llama-3.3-70b} : temps de réponse 1070ms, excellent français, coût compétitif
  \item \textbf{OpenAI GPT-4} : qualité supérieure mais coût 10x supérieur, latence variable
  \item \textbf{Anthropic Claude} : bonnes performances en français, restrictions d'usage
\end{itemize}

\subsection{Bases vectorielles : ChromaDB vs Cloud}

\begin{itemize}
  \item \textbf{ChromaDB} : open-source, intégration Python parfaite, contrôle total des données
  \item \textbf{Pinecone} : managé, scalabilité infinie mais dépendance cloud et coût mensuel
  \item \textbf{Weaviate} : fonctionnalités avancées mais complexité supérieure
\end{itemize}

\section{Performances et Métriques}

\begin{center}
\begin{tabularx}{0.95\textwidth}{X >{\centering\arraybackslash}C{3cm} >{\centering\arraybackslash}C{3cm}}
\toprule
\rowcolor{navyblue}\textcolor{white}{\textbf{Métrique}} & \textcolor{white}{\textbf{Objectif}} & \textcolor{white}{\textbf{Résultat}} \\
\midrule
\rowcolor{paleblue} Temps de réponse LLM & $<$ 2000ms & 1070ms moyen \\
Précision RAG & $>$ 85\% & 92\% (tests) \\
\rowcolor{paleblue} Validation complète & $<$ 500ms & 127ms moyen \\
Génération JSON & $<$ 300ms & 89ms moyen \\
\rowcolor{paleblue} Disponibilité API & 99.5\% & 99.8\% (30 jours) \\
\bottomrule
\end{tabularx}
\end{center}

\section{Conclusion}

L'Agent Hypothèses constitue une solution innovante qui transforme l'accompagnement entrepreneurial au Maroc en combinant intelligence artificielle, méthodologie financière rigoureuse et pédagogie adaptée. En collectant structurément les 22 hypothèses financières selon la méthode Chauvin, en les validant automatiquement et en générant trois scénarios prévisionnels, ce module IA démocratise l'accès à des outils financiers de qualité pour les entrepreneurs marocains. L'intégration d'un pipeline RAG avec 280 documents de référence et d'une API REST performante garantit la fiabilité et l'accessibilité de la solution, positionnant l'Agent Hypothèses comme un pont essentiel entre l'idée entrepreneuriale et la réalisation d'un business plan solide et validé.


\end{document}
